\section{Research Themes}
\label{sec:intro:themes}

This thesis addresses the gap between dynamics-aware and generalizable manipulation through three research themes, each targeting a different aspect of the problem.

\paragraph{RT1: Sampling-Based Kinodynamic Planning.}
The first theme develops planners that incorporate dynamics directly into the search process through physics simulation and kinodynamic exploration.
This includes introducing impulsive non-prehensile manipulation primitives, extending kinodynamic planning to settings with coupled robot-object dynamics, and developing multi-query planners that amortize dynamics computation across repeated queries.

\paragraph{RT2: Generative Models for Dynamics-Aware Motion.}
The second theme reframes motion generation as conditional sampling: a generative model learns to produce trajectories that implicitly satisfy dynamic constraints by conditioning on task-level properties such as payload mass or embodiment configuration.
This approach achieves constant-time generation without runtime constraint checking and generalizes across failure modes and payload levels without task-specific redesign.

\paragraph{RT3: Seeding for Constrained Trajectory Optimization.}
The third theme closes the loop between generative and optimization-based approaches by developing a task-agnostic trajectory seeding strategy.
By conditioning a diffusion model on robot state samples that satisfy dynamic constraints, this approach provides near-feasible initialization seeds that improve solver convergence and feasibility across diverse tasks without requiring task-specific engineering.

Together, these themes trace a progression from dynamics embedded in search, to dynamics embedded in learned representations, to dynamics embedded in initialization seeds for optimization.
